\chapter{Resources: The Context Layer}
\label{ch:resources}

\epigraph{``Roads? Where we're going, we don't need roads.''}{Doc Brown,
\emph{Back to the Future} (1985)}

Doc is wrong, of course. They need roads constantly for the rest of the
trilogy. But the line captures a real feeling, which is the moment you realise
the constraint you had been designing around is not actually the constraint.

With resources, the constraint people design around is ``how do I get data to
the model''. The actual constraint is ``how do I get data to the model
\emph{without} spending the context budget on data it will not read''. Those are
different problems, and confusing them is how you end up with a 200,000-token
context window and no room to think in.

\section{What a resource is, and what it is not}

A resource is addressable, readable context, identified by a URI.

The distinction from tools is about control, and it is worth being precise
because the specification is:

\begin{table}[htbp]
\centering
\small
\begin{tabularx}{\textwidth}{@{}lXX@{}}
\toprule
& \thead{Tools} & \thead{Resources} \\
\midrule
Controlled by & The model & The application \\
Invoked when  & The model decides & The host decides \\
Side effects  & Expected & None \\
Question      & ``Do something'' & ``Give me something to read'' \\
\bottomrule
\end{tabularx}
\caption{Model-controlled versus application-driven. The row that matters is the
first one.}
\label{tab:tools-vs-resources}
\end{table}

That first row is why resources exist at all. A user picking three documents
from a file tree is doing something a model should not be doing on their behalf,
and a host that wants to attach the current file to the conversation should not
have to persuade a language model to call a tool.

\begin{notebox}{When something could be either}
It happens constantly. ``Fetch the account summary'' is defensible as either.

The test I use: \emph{who should decide when this happens?} If the answer is
``the user, from a picker'', it is a resource. If the answer is ``the model,
while planning'', it is a tool. If the answer is genuinely both, ship both; they
can share an implementation, and Meridian's account summary is reachable as a
resource and as part of \texttt{assess\_account\_risk}.
\end{notebox}

\section{Mechanics}

Four methods, and they should look familiar by now.

\begin{wire}
resources/list             -> the concrete resources available
resources/templates/list   -> parameterised families of resources
resources/read             -> the contents of one URI
subscriptions/listen       -> tell me when these change
\end{wire}

A read looks like this:

\begin{wire}
{
  "jsonrpc": "2.0",
  "id": 2,
  "result": {
    "resultType": "complete",
    "contents": [
      {
        "uri": "meridian://accounts/ACC-1042/summary",
        "mimeType": "application/json",
        "text": "{\"accountId\":\"ACC-1042\",\"tier\":\"B\",...}"
      }
    ],
    "ttlMs": 900000,
    "cacheScope": "private"
  }
}
\end{wire}

Two rules worth internalising.

\textbf{A read may return several contents.} Reading a directory resource can
legitimately return every file in it. Useful, and a good way to accidentally put
four megabytes into a context window.

\textbf{An empty \texttt{contents} array is forbidden for a nonexistent
resource.} It is ambiguous: does the resource exist and have no content, or not
exist? Return \texttt{-32602} instead.

\begin{legacybox}{The error code changed}
Resource-not-found was \texttt{-32002} through \texttt{2025-11-25}. It is
\texttt{-32602} (Invalid params) now, aligning with plain JSON-RPC.

The asymmetry from Chapter~\ref{ch:transport} applies here specifically: read
the old code, write only the new one.
\end{legacybox}

\subsection{Templates}

Enumerating 240 accounts as concrete resources would be silly. A template
describes the family instead, using RFC 6570 URI templates:

\begin{py}
@server.template(
    "meridian://accounts/{account_id}/summary",
    "Account summary",
    description="Static profile for one account. Cheap, stable, private.",
    mime_type="application/json",
    ttl_ms=900_000,
    cache_scope="private",
)
def account_summary(ctx: RequestContext, uri: str, params: dict):
    account = ACCOUNTS.get(params["account_id"])
    if account is None:
        raise ResourceNotFound(uri)
    return json.dumps(account.to_json(), separators=(",", ":"))
\end{py}

Level 1 templates (a variable inside a path segment) cover essentially every
real case. Meridian compiles them to a regex once, at registration:

\begin{py}
def __post_init__(self) -> None:
    # RFC 6570 level 1 is all the book needs: {var} inside a path.
    pattern = re.escape(self.uri_template)
    pattern = re.sub(r"\\\{(\w+)\\\}", r"(?P<\1>[^/]+)", pattern)
    self._regex = re.compile("^" + pattern + "$")
\end{py}

\section{Designing URIs for machines and humans}

URIs are an interface, and unlike most interfaces this one is read by three
audiences: your code, the model, and a human staring at an audit log at
midnight. Design accordingly.

\textbf{Hierarchical and guessable.}
\texttt{meridian://accounts/ACC-1042/summary} tells you what it is without
documentation. \texttt{meridian://r/4a7f2b} does not.

\textbf{Stable.} A resource URI ends up in conversation history, in caches, in
audit logs, and in the model's working memory. Changing the scheme is a breaking
change even though nothing in the protocol says so.

\textbf{Scoped, so you cannot express a cross-tenant reference.}

\begin{dangerbox}{URIs that can name other people's data}
\texttt{meridian://accounts/\{id\}/summary} is fine only if the server checks
authorization on every read. If it does not, the URI space \emph{is} the attack:
a model that has seen one account id can construct another, and the identifiers
are sequential.

Two mitigations, and use both. Authorize every read against the caller. And
prefer identifiers that are not enumerable, so a leaked one does not imply the
neighbours.
\end{dangerbox}

\textbf{Pagination is a requirement, not a nicety.} Meridian's market-data
server has 40 benchmark resources and a page size of 25, which forces the
pagination path to be exercised rather than theoretical:

\begin{py}
def test_marketdata_paginates_its_resource_list(self):
    client = Client(InProcessTransport(marketdata.build_server()))
    first = client.call("resources/list")
    self.assertIn("nextCursor", first)
    everything = client.list_resources()
    self.assertGreater(len(everything), 25)
\end{py}

Cursors are opaque. Meridian uses an offset because its catalogue is small and
stable, and says so in the code, because an offset over a mutable table skips
and duplicates rows when the underlying set shifts between pages. If your list
comes from a database, encode a sort key instead.

\section{The 2026 caching model}

Here is the part of the chapter that will pay for the book.

Six operations carry caching hints: \texttt{server/discover},
\texttt{tools/list}, \texttt{prompts/list}, \texttt{resources/list},
\texttt{resources/templates/list}, and \texttt{resources/read}. Two fields:

\begin{description}[leftmargin=1.4em, style=nextline, itemsep=3pt]
\item[\texttt{ttlMs}] How long the client may consider this fresh. Semantically
identical to HTTP's \texttt{max-age}.
\item[\texttt{cacheScope}] \texttt{"public"} or \texttt{"private"}. Who may hold
the cached copy.
\end{description}

It is HTTP caching, and that is a compliment. HTTP caching is a design that has
survived thirty years of adversarial contact with reality.

\subsection{TTL is a freshness hint, not a polling interval}

The most common mistake, and it is a self-inflicted denial of service.

\keyidea{Check freshness when you need the data. Do not wake up every
\texttt{ttlMs} to refetch data nobody asked for. At scale, a fleet of clients
polling on a shared TTL is a synchronised thundering herd that you built
yourself, on purpose, out of enthusiasm.}

The correct loop is: need the data, check freshness, refetch only if stale.

\begin{py}
def get(self, server, method, params, auth_context="anon"):
    if method not in CACHEABLE_METHODS:
        self.stats.uncacheable += 1
        return None
    # A retry carrying MRTR fields is a different question with the same name.
    if "inputResponses" in params or "requestState" in params:
        self.stats.uncacheable += 1
        return None

    key = cache_key(server, method, params, auth_context)
    now = self._clock()
    with self._lock:
        entry = self._entries.get(key)
        if entry is None:
            self.stats.misses += 1
            return None
        if not entry.fresh_at(now):
            self.stats.stale += 1
            del self._entries[key]
            return None
        self.stats.hits += 1
        return entry.value
\end{py}

Note the second check. Results produced by an MRTR retry are never cacheable,
because they depend on \texttt{inputResponses}, which is not part of the cache
key. Caching one would serve one user's answer to another user's question.

Three edge cases the specification pins down, and Meridian tests all three:

\begin{itemize}[leftmargin=1.6em, itemsep=3pt]
\item \texttt{ttlMs: 0} means immediately stale. Do not store it.
\item \texttt{ttlMs} absent means assume 0. This should only happen with older
servers.
\item Negative \texttt{ttlMs} is ignored, treated as 0.
\end{itemize}

\subsection{\texttt{cacheScope} is an access-control decision}

\bookfig[0.98]{cache-scope}{The two scopes, and what each permits. The box at
the bottom is the part that turns a performance field into a security
field.}{cache-scope}

\texttt{public} means the response contains nothing user-specific, so any
client, gateway, or proxy may store it and serve it to anybody.

\texttt{private} means it may be reused only within the same authorization
context. Different token, different cache.

The trap is stated plainly in the specification and is worth repeating:
\textbf{a \texttt{public} response may be shared even when it came from an
authenticated endpoint}. Authentication does not make a response private.
Marking it \texttt{private} does.

So mislabelling a personalised response as \texttt{public} does not merely leak
it. It instructs a shared gateway to serve one user's data to another and report
a cache hit while doing so.

Meridian takes the conservative line and keys \emph{every} entry by principal,
public or not:

\begin{py}
def cache_key(server, method, params, auth_context):
    """Method plus the params that affect the answer, plus who is asking.

    `auth_context` is folded in unconditionally rather than only for `private`
    entries. Doing it the other way round means a single mislabelled response
    leaks across users, and the cost of being wrong is far higher than the cost
    of a few duplicate entries.
    """
\end{py}

That is a deliberate trade. It gives up gateway-level sharing of public
resources in exchange for making a mislabelling bug harmless. For a financial
application that is obviously right. For a public documentation server it may
not be, and you should make the choice consciously rather than inherit mine.

\subsection{What it is worth}

\bookfig[0.98]{cache-effect}{600 resource reads over 40 URIs with realistic
repetition. The left bar is what your server sees; the right column is what the
client experiences.}{cache-effect}

\begin{measurebox}{Honouring \texttt{ttlMs} on a realistic read pattern}
\begin{tabular}{@{}lrr@{}}
\toprule
\thead{} & \thead{Ignoring TTL} & \thead{Honouring TTL} \\
\midrule
Reads issued by the application & 600 & 600 \\
Requests reaching the server    & 600 & 23 \\
Wall clock                      & 12.4\,ms & 1.8\,ms \\
Cache hit rate                  & 0\,\% & 96.2\,\% \\
\midrule
\multicolumn{2}{@{}l}{Speedup} & 6.82$\times$ \\
\bottomrule
\end{tabular}

\vspace{6pt}
\footnotesize
In-process transport, so this measures the caching logic rather than the
network. Over a real network the ratio is far larger, because the 577 avoided
requests each avoid a round trip.
\end{measurebox}

Read the middle row again. 577 of 600 requests never happened. On loopback that
saves 10 milliseconds. Across a 68\,ms cross-region hop it saves 39 seconds.

\subsection{Invalidation}

TTL and change notifications are complementary, and the specification is explicit
that a server may provide either or both.

\begin{itemize}[leftmargin=1.6em, itemsep=3pt]
\item TTL alone: the client refetches when stale. Simple, and it may serve stale
data for up to \texttt{ttlMs}.
\item TTL plus \texttt{listChanged}: the notification invalidates immediately,
and the TTL stops needless refetching in between.
\end{itemize}

\begin{py}
def on_notification(self, msg: dict) -> None:
    """Turn a change notification into a cache invalidation.

    This is the whole reason `listChanged` and `ttlMs` coexist. The TTL
    stops you refetching while nothing has changed; the notification stops
    you serving stale data the moment something has.
    """
    method = msg.get("method", "")
    if method == "notifications/tools/list_changed":
        self.cache.invalidate_method(self.label, "tools/list")
    elif method == "notifications/resources/updated":
        uri = (msg.get("params") or {}).get("uri")
        if uri:
            self.cache.invalidate_uri(self.label, uri)
\end{py}

Subscribing is a \texttt{subscriptions/listen} request with a filter:

\begin{wire}
{
  "jsonrpc": "2.0",
  "id": 4,
  "method": "subscriptions/listen",
  "params": {
    "notifications": {
      "resourcesListChanged": true,
      "resourceSubscriptions": ["meridian://accounts/ACC-1042/summary"]
    }
  }
}
\end{wire}

Two rules govern that stream, and they exist because stdio multiplexes
everything onto one pipe:

\textbf{The server sends only what was requested.} Not ``everything it might
like''. A subscriber who asked for tool changes never sees resource updates.

\textbf{The acknowledgement comes first and carries the subscription id.}
Everything afterwards is tagged with it, so a client with two open subscriptions
can demultiplex.

\begin{py}
def test_acknowledged_first_then_only_requested_types(self):
    ...
    first = received[0]
    self.assertEqual(first["method"],
                     "notifications/subscriptions/acknowledged")
    self.assertEqual(
        first["params"]["_meta"]["io.modelcontextprotocol/subscriptionId"], 42)

    # A prompts change was never requested, so it must not be delivered.
    server.notify_list_changed("prompts")
    server.notify_list_changed("tools")
    ...
    self.assertIn("notifications/tools/list_changed", methods)
    self.assertNotIn("notifications/prompts/list_changed", methods)
\end{py}

The acknowledgement also reports the subset the server actually agreed to
honour, which is how a client learns that the server does not support resource
subscriptions rather than waiting forever for updates that are never coming.

\begin{legacybox}{\texttt{resources/subscribe} and the GET stream}
Through \texttt{2025-11-25}, subscribing was a \texttt{resources/subscribe} RPC
and notifications arrived on a standalone SSE stream opened with HTTP GET.

Both are gone, and for the same reason: they were per-connection state, which is
the thing this revision deleted. \texttt{subscriptions/listen} replaces them with
an ordinary request whose response happens to stay open.

On stdio, if the connection drops and is re-established, the client must re-send
\texttt{subscriptions/listen}. The server holds no subscription state across
reconnections.
\end{legacybox}

\section{Who chooses what goes into context}

Three strategies, three cost profiles, and most applications need all three.

\begin{table}[htbp]
\centering
\small
\begin{tabularx}{\textwidth}{@{}lXX@{}}
\toprule
\thead{Strategy} & \thead{How it works} & \thead{Token cost} \\
\midrule
User-driven
  & A picker, a file tree, an @-mention
  & Exactly what was chosen. Predictable, and usually right \\
Model-driven
  & Resource links in tool results; the model asks for what it wants
  & One extra round trip per fetch, but only fetches what it needs \\
Host policy
  & Always attach the current file, the system prompt, the last N turns
  & Constant per turn, and the easiest to let grow silently \\
\bottomrule
\end{tabularx}
\caption{Three ways context gets filled. The third one is where the waste
accumulates, because nobody has to decide anything for it to happen.}
\label{tab:context-strategies}
\end{table}

The model-driven route is worth a note. A tool can return a \emph{link} rather
than the content:

\begin{wire}
{
  "type": "resource_link",
  "uri": "meridian://filings/FIL-2000",
  "name": "FIL-2000",
  "description": "10-K, fiscal 2025, 284 pages",
  "mimeType": "application/json"
}
\end{wire}

The model sees that a 284-page filing exists and decides whether it wants it.
That is one extra round trip if it does, and zero tokens if it does not, which
is usually the better trade for large documents. Chapter~\ref{ch:agent-loop}
covers when to embed instead.

\section{Anti-patterns}

Five, in rough order of how often I encounter them.

\textbf{Dumping a database as resources.} One resource per row, ten thousand
rows. The list response alone exhausts the context window. If your resource list
is longer than a person would scroll, it should be a template plus a search
tool.

\textbf{Ignoring pagination.} Both directions. Servers that return everything,
and clients that read page one and stop. Meridian's client walks to the end and
caches each page under its own cursor.

\textbf{Cross-tenant URIs.} Covered above. Authorize on every read.

\textbf{\texttt{cacheScope: public} on anything personalised.} Covered above,
and it is a data breach rather than a performance bug.

\textbf{Returning enormous blobs with no size hint.} The \texttt{size} field
exists. Populate it. A host that knows a resource is 4\,MB before reading it can
decide not to, and a host that finds out afterwards has already spent the
budget.

\section{Meridian's resources}

\begin{table}[htbp]
\centering
\small
\begin{tabularx}{\textwidth}{@{}lXrl@{}}
\toprule
\thead{Resource} & \thead{What} & \thead{TTL} & \thead{Scope} \\
\midrule
\texttt{accounts/\{id\}/summary} & Account profile & 15\,min & private \\
\texttt{filings/\{id\}}          & Public regulatory filing & 24\,h & public \\
\texttt{risk/model-card}         & How the score works & 7\,d & public \\
\texttt{transactions/\{id\}}     & One transaction & 30\,s & private \\
\texttt{market/curve}            & Reference rates & 1\,h & public \\
\texttt{market/benchmark/\{...\}} & 40 sector benchmarks & 2\,h & public \\
\bottomrule
\end{tabularx}
\caption{Every TTL is a claim about how fast the underlying data changes, and
every scope is a claim about who may see it.}
\label{tab:meridian-resources}
\end{table}

The spread of TTLs is the interesting part, and each one is an argument:

\texttt{transactions} gets 30 seconds because a transaction under active review
changes state and stale data would mislead a compliance officer.

\texttt{filings} gets 24 hours and \texttt{public} because a filed 10-K is
immutable and identical for everybody. This is the row where a shared gateway
cache does real work.

\texttt{model-card} gets 7 days because it changes when the model is retrained,
which is quarterly.

\keyidea{A TTL is not a performance knob you tune until the graph looks nice. It
is a statement about how quickly the underlying data changes, and the right
value comes from the domain rather than from the dashboard.}

\section{What to remember}

\begin{itemize}[leftmargin=1.6em, itemsep=3pt]
\item Resources are application-driven; tools are model-controlled. The question
is who decides when it happens.
\item Templates over enumeration. Pagination is required, and cursors are opaque.
\item Six operations carry \texttt{ttlMs} and \texttt{cacheScope}.
\item TTL is a freshness hint. Never poll on it.
\item MRTR retries and interim results are never cacheable.
\item \texttt{cacheScope: public} may be shared across users even from an
authenticated endpoint. It is an access-control decision.
\item Honouring TTL avoided 577 of 600 requests on a realistic read pattern, a
96.2\,\% hit rate.
\item Notifications invalidate; TTL prevents needless refetching between them.
The server sends only the types the client asked for.
\item Authorize every read. Populate \texttt{size}. Do not dump your database.
\end{itemize}

Next: prompts, the primitive everybody skips, and the one that quietly
standardises a workflow across three teams.
